\chapter{Conclusion}
\label{cha:conclusion}

This thesis investigated the latency behavior of agentic LLM serving under position-independent KV cache reuse, delivering a set of backend enhancements to LMCache, the reference implementation of CacheBlend, together with an evaluation that localizes the bottleneck that remains. The starting observation was that agentic prompts differ from the multi-round QA workloads that CacheBlend was originally evaluated on: they contain code blocks that are reordered or partially substituted across successive requests, arrive at higher concurrency, and span a far wider range of cache hit ratios. The upstream implementation leaves four aspects of this regime unaddressed, and this work delivered a corresponding set of fixes.

\section{Summary of Contributions}
\label{sec:conclusion_contributions}

The four backend enhancements are as follows. First, we added the \texttt{old\_positions} field to \texttt{MemoryObjMetadata} together with its serialization handlers, populated it at save time, and connected it to the layerwise GPU connector's rotary embedding correction path. This completes the runtime handling for reusing a cached chunk at a position different from where it was originally cached, which is a common pattern in agentic workflows. Second, we introduced a \texttt{batched\_allocate\_varied} method in the CPU storage backend that acquires the buffer lock once per request rather than once per chunk, reducing lock contention proportionally to the number of chunks. Third, we bounded the eviction loop with a configurable retry limit (\texttt{MAX\_EVICTION\_RETRIES}), which prevents system-wide stalls when all cache entries are pinned by concurrent lookups. Fourth, we implemented a per-request profiling module that instruments the critical path with per-operation and per-phase timing records, enabling fine-grained latency attribution that would otherwise be invisible in aggregate statistics.

These enhancements were integrated into an end-to-end system in which the RepoGraph + Agentless frontend generates realistic agentic prompts, an adapter module bridges the frontend prompt format to the backend cache engine, and the backend serves the requests with CacheBlend enabled. A request logging and replay mechanism was added so that the same trace can drive the backend under both blend and no-blend configurations, isolating the backend effect from agent-side variability.

\section{Summary of Findings}
\label{sec:conclusion_findings}

The latency evaluation on a 537-request SWE-bench Verified-mini trace produced four main findings.

First, cache hit ratio, not prompt length, determines whether blending pays off. The trace splits into a low-hit cluster (hit ratio at most 34.2\%, $n = 430$, median speedup 0.496) in which 2.1\% of requests beat the baseline, and a high-hit cluster (hit ratio at least 72.3\%, $n = 106$, median speedup 2.058) in which 81.1\% do. No request in the trace falls between these two hit-ratio bands, so the breakeven crossing can be bounded but not measured.

Second, applied to all traffic, the cache layer is a net cost on this workload. Only 17.7\% of requests were faster under blending, the median speedup is 0.509$\times$, and summed over the trace the blend configuration spends 772\,s more than vanilla vLLM would have.

Third, the reason is that the KV store runs synchronously inside the model execution step, before the first token is produced, and transfers a volume linear in prompt length. A request therefore pays for what it writes as well as what it reads, and because that cost grows in proportion to the baseline, the penalty for a cache miss stays near 0.5$\times$ at every prompt length from under 1{,}000 to over 12{,}000 tokens. Prompt length does not rescue it.

Fourth, short requests exhibit far higher speedup variance than long ones. This variance is driven by an interaction between vLLM's chunked-prefill scheduler and LMCache's per-step invocation pattern: requests that arrive during decode-heavy scheduling windows are fragmented across many prefill steps, and each step pays the per-call retrieve overhead of 30.8\,ms in blend mode versus 18.0\,ms without blending, of which GPU onload is 91\%. This mechanism explains the spread among short requests but not the aggregate deficit, which the store path accounts for.

The frontend evaluation, run as a supporting cross-check, confirmed two properties of the integrated system. RepoGraph produces comparable resolve rates across four LLM backends (GPT-4o, o4-mini, Claude 3.5 Sonnet, Devstral Small), which supports its portability claim. Enabling CacheBlend at various blending granularities on Devstral Small yields resolve rates comparable to the no-blend baseline, which confirms that the position correction enhancement preserves output quality end to end.

\section{Limitations and Future Work}
\label{sec:conclusion_future}

Several limitations of the current evaluation suggest concrete next steps. The most important is that the SWE-bench Verified-mini trace samples cache hit ratio very unevenly. Most requests share little content with their predecessors, and no request at all falls in the 34.2\% to 72.3\% hit-ratio band, which is precisely where the breakeven crossing lies. A workload constructed to sweep hit ratio continuously would locate that crossing and turn the bound reported here into a measurement. This matters directly for the mitigation strategies below, since a bypass policy needs to know where the crossing is. The bimodal fragmentation observed in the trace is best interpreted as an emergent property of steady-state scheduler load rather than as a hardcoded system constant. Running the same benchmark under different concurrency levels would characterize how the mode shifts and would generalize the fragmentation discussion.

The evaluation also points to a specific implementation change that was not made here. Section~\ref{sec:rec_miss_penalty} traces the roughly 2$\times$ TTFT penalty on cache misses to the fact that the KV store runs synchronously inside the model execution step, before the first token is produced, and transfers a volume linear in prompt length. Because roughly 80\% of the requests in our trace miss, this single property accounts for most of the aggregate slowdown we measured. Moving the offload off the critical path, so that a request pays for what it reads but not for what it writes, is the highest-value next step this work identifies. Alongside it, three mitigations for the small-request variance were identified but not implemented: a minimum token-per-step threshold for micro-batched steps, the prompt-length bypass evaluated in Section~\ref{sec:rec_length_floor}, and a coalesced retrieval mode that batches token ranges across prefill steps.

Beyond these targeted mitigations, the frontend interface is intentionally narrow, and swapping RepoGraph + Agentless for a more recent context retrieval module (for example, one based on agentic search rather than static graph queries) would provide a useful comparison point and would validate the pluggable frontend claim in practice.

Taken together, the results show that position-independent reuse is profitable for agentic workloads in the high-hit-ratio regime, but that applying it to all traffic is a net cost until the KV store is moved off the critical path. Prompt length turns out to be a poor proxy for that regime: it separates the profitable requests from the unprofitable ones only weakly, and the length-based routing rules we evaluated reduce the aggregate cost without eliminating it. The mechanism itself is sound, and the enhancements delivered here make it correct and diagnosable under agentic conditions. What the evaluation adds is a specific account of where the remaining cost lives and what it would take to remove it, which is the necessary step between a promising QA-workload benchmark and a production-ready component in a code agent.
